\chapter{Grant Semantics}\label{chap:grants}

The core evaluation uses deny-overrides semantics: a request is allowed
iff no matching Deny statement applies and some matching Allow statement
applies.  This chapter presents the generic deny-overrides theory
(vendor-neutral, in \texttt{Seclib.Domain.PolicySem}) and its
instantiation to AWS IAM (in \texttt{IamExplainer.Proofs}).

\section{Generic Deny-Overrides}

\begin{definition}[denyOverrides]\label{def:denyOverrides}
\lean{Domain.denyOverrides}
\leanok
Generic deny-overrides evaluation parameterized over a \texttt{PolicySem}
instance.  Returns \texttt{true} iff no deny applies ($\texttt{condEval} = T$)
and some allow applies ($\texttt{condEval} \neq F$).
\end{definition}

\begin{theorem}[removeAllow]\label{thm:removeAllow}
\lean{Domain.removeAllow}
\leanok
Removing an Allow statement from a list can only narrow the set of
allowed requests.  If \texttt{denyOverrides} on
\texttt{stmts.eraseIdx i} allows a request, so does \texttt{denyOverrides}
on the original \texttt{stmts}.
\end{theorem}

\begin{theorem}[narrow\_narrows]\label{thm:narrow_narrows}
\lean{Domain.narrow_narrows}
\leanok
Replacing an Allow statement with one that matches fewer requests
(same effect, same conditions, monotone \texttt{sat}) narrows access.
\end{theorem}

\begin{theorem}[filterMap\_narrows]\label{thm:filterMap_narrows}
\lean{Domain.filterMap_narrows}
\leanok
Applying a \texttt{filterMap} that preserves Deny statements and
produces effect-preserving, condition-preserving, match-monotone
replacements for Allow statements narrows access.
\end{theorem}

\begin{theorem}[nocontext\_conservative]\label{thm:nocontext_conservative}
\lean{Domain.nocontext_conservative}
\leanok
If allowed under any context, then allowed under \texttt{noCtx}.
Relies on the \texttt{PolicySem} axioms \texttt{cond\_noCtx\_tu}
and \texttt{cond\_noCtx\_T\_imp}.
\end{theorem}

\begin{theorem}[grants\_complete (generic)]\label{thm:grants_complete_generic}
\lean{Domain.grants_complete}
\leanok
If \texttt{denyOverrides} allows a request, there exists an Allow
statement in the list that matches and whose condition is not~F.
\end{theorem}

\section{Concrete Rule-Level Theory}

\begin{definition}[appliesOf]\label{def:appliesOf}
\lean{appliesOf}
\leanok
Determines whether a rule applies given its effect and the
three-valued condition result: Allow applies when $\neq F$,
Deny applies when $= T$.
\end{definition}

\begin{definition}[denyOverridesRules]\label{def:denyOverridesRules}
\lean{denyOverridesRules}
\leanok
Concrete deny-overrides on a list of \texttt{Rule} values.
\end{definition}

\begin{theorem}[denyOverrides\_remove\_allow\_narrows]\label{thm:denyOverrides_remove_allow_narrows}
\lean{denyOverrides_remove_allow_narrows}
\leanok
Removing an Allow rule narrows access (concrete rule-level).
\end{theorem}

\begin{theorem}[denyOverrides\_add\_deny\_narrows]\label{thm:denyOverrides_add_deny_narrows}
\lean{denyOverrides_add_deny_narrows}
\leanok
Adding a Deny rule narrows access (concrete rule-level).
\end{theorem}

\begin{lemma}[conj\_narrows]\label{lem:conj_narrows}
\lean{conj_narrows}
\leanok
Boolean conjunction: if $a \land b = \text{true}$ then both
$a = \text{true}$ and $b = \text{true}$.
\end{lemma}

\begin{theorem}[appliesOf\_unknown\_conservative]\label{thm:appliesOf_unknown_conservative}
\lean{appliesOf_unknown_conservative}
\leanok
When conditions are unknown (T or U), the rule-level evaluation
over-approximates: the no-context result is at least as permissive.
\end{theorem}

\section{IAM Instantiation}

\begin{definition}[stmtGrantsAction]\label{def:stmtGrantsAction}
\lean{stmtGrantsAction}
\leanok
Tests whether a statement grants a specific action string, checking
\texttt{actions} and \texttt{notActions} with case-insensitive glob matching.
\end{definition}

\begin{definition}[allows]\label{def:allows}
\lean{allows}
\leanok
Policy-level evaluation: deny-overrides over the policy's statements.
\end{definition}

\begin{definition}[iamSem]\label{def:iamSem}
\lean{iamSem}
\leanok
The \texttt{PolicySem} instance for AWS IAM, wiring \texttt{Statement.effect},
\texttt{stmtMatches}, \texttt{evalCond}, and \texttt{noContext}.
\end{definition}

\begin{corollary}[removeAllow\_narrows]\label{cor:removeAllow_narrows}
\uses{thm:removeAllow}
\lean{removeAllow_narrows}
\leanok
Removing an Allow statement from an IAM policy narrows access.
Direct application of the generic \texttt{removeAllow}.
\end{corollary}

\begin{theorem}[matchPattern\_ci\_congr]\label{thm:matchPattern_ci_congr}
\uses{lem:ciEq_toLowerStr}
\lean{matchPattern_ci_congr}
\leanok
\texttt{matchActionPattern} is congruent under \texttt{ciEq}:
if the pattern matches \texttt{s}, it matches any \texttt{s'}
with \texttt{ciEq s s' = true}.
\end{theorem}

\begin{theorem}[stmtGrantsAction\_ci\_congr]\label{thm:stmtGrantsAction_ci_congr}
\uses{lem:ciEq_toLowerStr}
\lean{stmtGrantsAction_ci_congr}
\leanok
\texttt{stmtGrantsAction} is congruent under \texttt{ciEq}.
\end{theorem}

\begin{theorem}[stmtGrantsAction\_narrow]\label{thm:stmtGrantsAction_narrow}
\uses{lem:matchPatternGo_literal_eq, lem:toLower_eq_star, lem:toLower_eq_question, thm:stmtGrantsAction_ci_congr}
\lean{stmtGrantsAction_narrow}
\leanok
If \texttt{new\_stmt} has explicit wildcard-free actions each granted
by the original statement, and \texttt{stmtGrantsAction new\_stmt a = true},
then \texttt{stmtGrantsAction s a = true}.
\end{theorem}

\begin{theorem}[allows\_replaceAllow\_mono]\label{thm:allows_replaceAllow_mono}
\uses{thm:narrow_narrows}
\lean{allows_replaceAllow_mono}
\leanok
Generic Allow-statement replacement monotonicity for IAM policies.
If the replacement preserves resources, notResources, condition, and is
grant-monotone on actions, then the replacement narrows \texttt{allows}.
\end{theorem}

\begin{theorem}[narrowActions\_narrows]\label{thm:narrowActions_narrows}
\uses{thm:allows_replaceAllow_mono, thm:stmtGrantsAction_narrow}
\lean{narrowActions_narrows}
\leanok
\textbf{Gate theorem (T1).}
Action narrowing narrows the policy.  Hypotheses: the replacement
statement is Allow, uses \texttt{actions} (not \texttt{notActions}),
every action string is a wildcard-free literal that the original
statement already grants, resources/notResources/condition are
preserved.  Conclusion: every request allowed by the narrowed policy
was already allowed by the original.
\end{theorem}

\begin{theorem}[allows\_replace\_stmt\_mono]\label{thm:allows_replace_stmt_mono}
\uses{thm:narrow_narrows}
\lean{allows_replace_stmt_mono}
\leanok
Generalized statement replacement monotonicity: if the replacement
preserves effect (Allow), condition, and is match-monotone, then
replacing narrows \texttt{allows}.
\end{theorem}

\begin{theorem}[narrowResources\_narrows]\label{thm:narrowResources_narrows}
\uses{thm:allows_replace_stmt_mono, thm:matchPattern_cs_literal_eq}
\lean{narrowResources_narrows}
\leanok
\textbf{Transform T3.}
Replacing resource patterns with wildcard-free literals each matched
by the old patterns narrows \texttt{allows}.  Hypotheses include
\texttt{hnew\_notres = none} and each replacement resource is
wildcard-free and matched by the original.
\end{theorem}
